\worldchapter{Spielleitung und Gegner}{pantheon}
\label{ch:spielleitung}

\begin{multicols*}{2}
\section{Gegner in unter einer Minute}
Das Spiel will keine langen Monsterbau-Kapitel. Fast jede Gegenseite lasst sich aus wenigen Werten zusammensetzen.

\section{Schneller Statblock}
Jeder NSC oder jedes Wesen braucht nur:
\begin{itemize}[leftmargin=*]
\item Name und Rolle
\item drei Zielwerte: \textbf{Assault}, \textbf{Reaction}, \textbf{Will}
\item Schaden 1--3
\item Schutz 0--3
\item Wound Threshold 2--8
\item einen Signature Move
\item eine ausnutzbare Schwache
\end{itemize}

\section{Standard-Zielwerte nach Tier}
\begin{description}[style=nextline,leftmargin=0pt]
\item[\textbf{Frail}] 40
\item[\textbf{Trained}] 50
\item[\textbf{Veteran}] 60
\item[\textbf{Elite}] 70
\item[\textbf{Monstrous}] 80
\end{description}

Wahle ein Tier fur alle drei Zielwerte und verschiebe dann einen Zielwert um +10 und einen anderen um -10, um den Gegner zu akzentuieren.

\textbf{Reaction} ist immer der Wert fur Parieren, Ausweichen oder sonstige Abwehr gegen eingehende Angriffe.

\section{Schaden und Schutz}
\begin{itemize}[leftmargin=*]
\item \textbf{Schaden 1}: zivile Waffe, schwacher Naturangriff
\item \textbf{Schaden 2}: Kriegswaffe, Raubtierhieb
\item \textbf{Schaden 3}: schwere Waffe, Ogerwucht, groBer Zauber
\end{itemize}

\begin{itemize}[leftmargin=*]
\item \textbf{Schutz 0}: ungeschutzt
\item \textbf{Schutz 1}: gepolstert, Fell, leichte Haut
\item \textbf{Schutz 2}: Kette, Platte, verharter Panzer
\item \textbf{Schutz 3}: Reliktpanzer oder Monsterschale
\end{itemize}

\section{Wound Threshold}
\begin{itemize}[leftmargin=*]
\item \textbf{Minion}: 2--3
\item \textbf{Standardgegner}: 4--5
\item \textbf{Elite}: 6--7
\item \textbf{Boss}: 8+
\end{itemize}

\section{Gruppengegner}
Schwarme, Rotten oder Mobgegner konnen als gemeinsamer Block laufen.
\begin{itemize}[leftmargin=*]
\item jede \textbf{Serious Consequence} entfernt 1 Einheit
\item bei halber Einheitenzahl bekommt der Mob -10 auf Offensive
\item bei 0 Einheiten ist die Gruppe zerschlagen oder geflohen
\end{itemize}

\section{Signature Moves}
Ein guter Signature Move tut genau eines:
\begin{itemize}[leftmargin=*]
\item ein Ziel versetzen
\item einen Zustand auferlegen
\item die Gruppe trennen
\item eine harte Wahl oder einen unmittelbaren Preis erzwingen
\end{itemize}

Er sollte als ein Satz formulierbar sein und mit genau einem Wurf funktionieren.

\section{Schwachen}
Eine Schwache muss konkret ausnutzbar sein, etwa:
\begin{itemize}[leftmargin=*]
\item bestimmter Schadenstyp
\item Ritualwort
\item Terrainabhangigkeit
\item blinder Sinnesfleck
\end{itemize}

Wird die Schwache richtig ausgespielt, sinkt fur diesen Austausch ein Zielwert des Gegners um -20.

\begin{ruleexample}
\textbf{Grubenhund der Krypta}\\
Trained, Assault 60, Reaction 50, Will 40, Schaden 2, Schutz 1, Threshold 4.\\
Signature Move: ``Zerrt bei Treffer das Ziel aus Near nach Engaged und macht es Shaken.''\\
Schwache: ``Blendendes geweihtes Licht senkt Reaction fur einen Austausch um 20.'' 
\end{ruleexample}
\end{multicols*}
